\section*{Bab \ref{ch:g12:continental-crust} --- Kerak Benua}

\begin{solution}{exo:g12:continental-crust:1}
\omterm{def:g12:continental-crust:continental-crust}{Kerak benua}: 30 sampai \qty{70}{km}, rata-rata sekitar \qty{41}{km}; batuan
granitik, gneis, dan batuan metamorf serta \omterm{def:g11:igneous-rocks:volcanic-plutonic}{plutonik} lain di bawah selimut
sedimen; rata-rata sekitar \qty{61}{\percent} silika; sekitar
\qty{2830}{kg/m^3}; sampai 4.03 miliar tahun (butir-butir 4.4). \omterm{def:g11:ocean-ridges:oceanic-lithosphere}{Kerak samudra}:
sekitar \qty{7}{km}; basalt dan \omterm{def:g11:ocean-ridges:gabbro}{gabro}; sekitar \qty{50}{\percent} silika;
sekitar \qty{3000}{kg/m^3}; kurang dari 200 juta tahun.
\end{solution}

\begin{solution}{exo:g12:continental-crust:2}
\qty{60.6}{\percent}: \omterm{def:g11:igneous-rocks:felsic-mafic}{intermediet}, seperti \omterm{def:g11:igneous-rocks:rock-names}{diorit}. \qty{66.6}{\percent}:
\omterm{def:g11:igneous-rocks:felsic-mafic}{felsik}, seperti granit.
\end{solution}

\begin{solution}{exo:g12:continental-crust:3}
\omterm{def:g11:ocean-ridges:oceanic-lithosphere}{Kerak samudra}, sekitar \qty{3000}{kg/m^3}, terletak di atas mantel \omterm{def:g11:lithosphere-asthenosphere:lithosphere}{litosfer},
dan \omterm{def:g11:ocean-ridges:oceanic-lithosphere}{litosfer samudra} yang tua dan dingin secara keseluruhan cukup rapat untuk
tenggelam. \omterm{def:g12:continental-crust:continental-crust}{Kerak benua}, sekitar \qty{2830}{kg/m^3}, jauh lebih ringan daripada
mantel (sekitar \qty{3380}{kg/m^3}), terlalu besar daya apungnya: benua yang
tiba di palung menghentikan \omterm{def:g11:subduction-zones:subduction}{subduksi}.
\end{solution}

\begin{solution}{exo:g12:continental-crust:4}
(1) Inti benua yang tua dan stabil, yang tidak lagi terdeformasi. (2) Lebih
dari satu miliar tahun. (3) Batuan metamorf dan \omterm{def:g11:igneous-rocks:volcanic-plutonic}{plutonik}, \omterm{def:g12:continental-crust:root}{akar-akar} dalam
pegunungan masa lalu.
\end{solution}

\begin{solution}{exo:g12:continental-crust:5}
$4030/190 \approx 21$ kali. Dasar samudra dibuat di punggungan dan
dihancurkan di \omterm{def:g11:subduction-zones:subduction}{zona subduksi}: tidak ada bagiannya yang bertahan jauh lebih lama
daripada 200 juta tahun.
\end{solution}

\begin{solution}{exo:g12:continental-crust:6}
(1) $0.306 \times 210.4 \approx \num{64.4}\text{ juta }\unit{km^2}$.
(2) $210.4 - 64.4 = \num{146.0}\text{ juta }\unit{km^2}$.
(3) $4\pi \times 6371^2 \approx \qty{5.10e8}{km^2}$; $210.4/510.1 \approx
0.41$: \qty{41}{\percent}.
\end{solution}

\begin{solution}{exo:g12:continental-crust:7}
(1) $41100 \times 2830 \approx \qty{1.16e8}{kg}$. (2) $7100 \times 3000 =
\qty{2.13e7}{kg}$. (3) Sekitar 5.5 kali.
\end{solution}

\begin{solution}{exo:g12:continental-crust:8}
\omterm{def:g12:continental-crust:root}{Akar}: $2.5 \times 2830/550 \approx \qty{12.9}{km}$. Keraknya lebih tebal
sebesar $2.5 \times 3380/550 \approx \qty{15.4}{km}$: seluruhnya sekitar
\qty{50}{km}, dalam rentang 50 sampai \qty{70}{km} yang ditemukan di bawah
sabuk pegunungan.
\end{solution}

\begin{solution}{exo:g12:continental-crust:9}
(1) Sekitar $0.103 + 0.100 + 0.102 \approx 0.30$: \qty{30}{\percent}.
(2) Pita yang termuda, 0 sampai 10 juta tahun (sekitar \qty{10}{\percent}),
hampir sama dengan pita 20 sampai 30. (3) Dasar yang tua telah dihancurkan di
\omterm{def:g11:subduction-zones:subduction}{zona subduksi}; makin tua dasarnya, makin banyak bagiannya yang telah hilang.
\end{solution}

\begin{solution}{exo:g12:continental-crust:10}
$120 = 0.694\,h - 0.306 \times 100$, sehingga $h = 150.6/0.694 \approx
\qty{217}{m}$. Bila hanya daratan yang dihitung, tepi-tepi benua yang rendah
dan terendam tidak diikutkan, dan rata-ratanya menjadi jauh lebih tinggi.
\end{solution}

\begin{solution}{exo:g12:continental-crust:11}
Pulau itu: \omterm{def:g11:ocean-ridges:gabbro}{gabro} di bawah lava dan \omterm{def:g11:seismic-structure:moho}{Moho} pada \qty{8}{km}, \omterm{def:g11:ocean-ridges:oceanic-lithosphere}{kerak samudra}
(ketinggian pulau itu tidak mengubah hal tersebut). Dataran itu: gneis dan
\omterm{def:g11:seismic-structure:moho}{Moho} pada \qty{38}{km}, \omterm{def:g12:continental-crust:continental-crust}{kerak benua}.
\end{solution}

\begin{solution}{exo:g12:continental-crust:12}
(1) Eksak: $\tfrac{4}{3}\pi(6371^3 - 6329.9^3) \approx
\qty{2.0829e10}{km^3}$; pendekatan: $4\pi \times 6371^2 \times 41.1
\approx \qty{2.0964e10}{km^3}$. (2) Pendekatan itu terlalu besar sekitar
\qty{0.65}{\percent}, dekat dengan $H/R \approx 0.0065$: untuk lapisan yang
tipis, galatnya dapat diabaikan.
\end{solution}

\begin{solution}{exo:g12:continental-crust:13}
(1) $60.6 = (66.6 + 2x)/3$, sehingga $x = (181.8 - 66.6)/2 = \qty{57.6}{\percent}$.
(2) \omterm{def:g11:igneous-rocks:felsic-mafic}{Intermediet}. (3) \omterm{def:g12:continental-crust:continental-crust}{Kerak benua} bagian dalam kurang kaya silika daripada kerak
bagian atas: lebih banyak mineral \omterm{def:g11:igneous-rocks:felsic-mafic}{mafik}, lebih kaya besi, magnesium, dan
kalsium.
\end{solution}

\begin{solution}{exo:g12:continental-crust:14}
(1) $3 \times 10^8/3 = 10^8$ tahun, 100 juta tahun. (2) Orde besarnya sama
dengan umur dasar samudra: tiga perempatnya lebih muda daripada 100 juta tahun,
tidak ada yang lebih tua daripada sekitar 190. (3) Dasar itu pasti dihancurkan
secepat dasar itu dibuat: dasar itu tenggelam kembali ke dalam mantel di \omterm{def:g11:subduction-zones:subduction}{zona
subduksi}.
\end{solution}

\begin{solution}{exo:g12:continental-crust:15}
(1) $8.6 \times 10^9 / 4.0 \times 10^9 \approx \qty{2.2}{km^3}$ setahun.
(2) Laju yang tetap hanyalah satu kemungkinan: kerak mungkin tumbuh cepat pada
masa awal lalu didaur ulang sejak itu, atau tumbuh belakangan. Batuan dan
butir-butir tertua membuktikan bahwa sebagian kerak sudah ada sejak awal;
keduanya tidak dapat mengatakan seberapa banyak, karena kerak tua dapat
tererosi atau terbawa kembali ke dalam mantel dan hanya meninggalkan sedikit
jejak.
\end{solution}

\begin{solution}{pb:g12:continental-crust:1}
\textbf{1.} $4\pi \times 6371^2 \approx \qty{5.10e8}{km^2}$.

\textbf{2.} $210.4/510.1 \approx 0.41$: \qty{41}{\percent}.

\textbf{3.} $0.306 \times 210.4 \approx \num{64.4}\text{ juta }\unit{km^2}$.

\textbf{4.} $510.1 - 210.4 \approx \num{299.7}\text{ juta }\unit{km^2}$,
\qty{59}{\percent}.

\textbf{5.} Tepi-tepi benua terletak di bawah laut dangkal, pada paparan
benua.

\textbf{6.} Tebalnya, \qty{41}{km}, sangat kecil dibandingkan dengan
jari-jarinya, \qty{6371}{km}: kelengkungan mengubah volumenya kurang dari satu
persen.

\textbf{7.} $210.4 \times 10^6 \times 41.1 \approx \qty{8.65e9}{km^3}$.

\textbf{8.} \qty{8.65e18}{m^3}.

\textbf{9.} $210.4 \times 10^6 \times 36 \approx \qty{7.57e9}{km^3}$.

\textbf{10.} $(8.65 - 7.57)/7.57 \approx 0.14$: \qty{14}{\percent}.

\textbf{11.} $8.65 \times 10^{18} \times 2830 \approx \qty{2.45e22}{kg}$.

\textbf{12.} $41100 \times 2830 \approx \qty{1.16e8}{kg}$.

\textbf{13.} $299.7 \times 10^6 \times 7.1 \approx \qty{2.13e9}{km^3}$, yaitu
\qty{2.13e18}{m^3}; dikali 3000: sekitar \qty{6.4e21}{kg}.

\textbf{14.} \omterm{def:g12:continental-crust:continental-crust}{Kerak benua}, sekitar 3.8 kali lebih berat: kerak itu hampir enam
kali lebih tebal, dan hal ini melampaui pengaruh luas dan massa jenisnya yang
lebih kecil.

\textbf{15.} \omterm{def:g12:continental-crust:continental-crust}{Kerak benua} tersusun dari batuan granitik, yang kaya kuarsa dan
feldspar, mineral pembentuk batuan yang paling ringan; \omterm{def:g11:ocean-ridges:oceanic-lithosphere}{kerak samudra} tersusun
dari basalt dan \omterm{def:g11:ocean-ridges:gabbro}{gabro}, yang kaya piroksen dan olivin.

\textbf{16.} $2.45 \times 10^{22}/5.97 \times 10^{24} \approx 0.0041$.

\textbf{17.} $3 \times 0.41 \times (41.1/6371) \times (2830/5513) \approx
0.0041$: sama.

\textbf{18.} $6.4 \times 10^{21}/5.97 \times 10^{24} \approx 0.0011$: sekitar
sepersepuluh persen.

\textbf{19.} Sekitar \qty{0.4}{\percent}.

\textbf{20.} \omterm{def:g12:continental-crust:continental-crust}{Kerak benua} adalah sekitar \textbf{\qty{0.4}{\percent} dari massa
Bumi}.
\end{solution}
